\section{Introduction}
Privacy amplification is key in differentially private (DP) machine learning (ML) as it enables tighter privacy budgets under certain assumptions on the data processing.
For example, one of the main contributions in the DP-SGD (DP Stochastic Gradient Descent) work by~\citet{abadi2016deep} was the ``moments accountant'', which relies on privacy amplification~\citep{KLNRS,BST14} for bounding the privacy cost. Recently, privacy amplification analysis enabled~\citet{BandedMF} to show that a class of DP-FTRL (DP Follow-The-Regularized-Leader) algorithms~\citep{thakurta2013nearly,kairouz2021practical, MEMF} is superior in privacy-utility tradeoffs to DP-SGD.\footnote{Precisely, they showed DP-FTRL is never worse, and often better, than DP-SGD---it ``pareto-dominates''.}
At the heart of DP-FTRL is the matrix mechanism~\citep{mckenna2021hdmm,denisovDPMF}. Thus, bringing privacy amplification to matrix mechanisms (\mms) is an important area of research to enable better privacy-utility tradeoffs.

The \mm effectively computes the prefix sums $\sum_{i\leq t} \bfx_i$ over a sequence of adaptively chosen vectors $\{\bfx_i:i\in[n]\}$. This can be written as computing $\bfA\cdot\bfx$ where $\bfA$ is a lower triangular matrix with all ones and $\bfx=[\bfx_1|\cdots|\bfx_n]^\top\in\calR^{n\times d}$. Observe that
releasing $\bfA\bfx[i]$ releases the model parameters at step $i$ in SGD and that if $\bfx$ is the clipped and averaged model gradient (e.g., as returned by any ML optimizer like SGD), then the DP release of $\bfA\bfx$ gives a DP-FTRL optimizer.
%
The matrix mechanism factorizes $\bfA=\bfB\cdot\bfC$ to minimize the error (introduced in $\bfB$) in the prefix sum estimates, while ensuring $\bfC\cdot \bfx +\bfz$ satisfies DP, where $\bfz$ is drawn from an isotropic normal distribution. We refer to $\bfB$ as the decoder and $\bfC$ as the encoder.

\mms pose a major challenge for privacy amplification analysis. Standard privacy amplification exploits randomness in the selection of minibatches\footnote{E.g., for a data set $D$, a row $\bfx_{[i,:]}=\sum_{d\in S}\nabla_\theta \ell(\theta_i;d)$, where $S$ is a randomly chosen subset of $D$ (e.g., sampled uniformly at random from $D$, or a subset from a random shuffling of $D$), $\ell$ is a loss function, and $\theta_i$ is obtained via an SGD state update process.} but requires that the noise added to each minibatch is independent. In the matrix mechanism, a minibatch ($\bfx_i$) contributes to multiple rows of $\bfC\cdot\bfx + \bfz$, leading to correlated noise that prevents direct application of amplification. This challenge can be seen by the limitations of the amplification analysis of~\citet{BandedMF} which only applies to a special class of `$b$-banded' matrix mechanisms (i.e., the first $b$ principal diagonals of $\bfC$ are non-zero), that in-turn leads to multiplicatively higher sampling probabilities
preventing the full benefits of amplification.
Resulting from these limitations is a large range of $\epsilon$ where banded matrix mechanisms cannot simultaneously leverage the benefits of correlated noise and privacy amplification; in other words, they perform equivalent to, \textit{but no better than, DP-SGD}\footnote{In~\citet{BandedMF}, this region surfaces empirically even for larger $\epsilon\approx1$.}.  Further, since their analysis only applies to the special banded case, matrix mechanisms from the extant literature cannot leverage amplification and correlated noise, e.g.,~\cite{kairouz2021practical,MEMF,denisovDPMF}. 

\begin{center}
\textit{In this work, we provide a generic privacy amplification machinery for adaptive matrix mechanisms for any lower-triangular encoder matrix $\bfC$ that strictly generalizes the approach in~\citep{BandedMF}.}
\end{center}
%
\vspace{-1.5em}  % SPACE
\subsection{Our contributions}
\vspace{-0.25em}  % SPACE
Our main contribution is to prove a general privacy amplification analysis for \textit{any matrix mechanism}, i.e., arbitrary encoder matrices $\bfC$ for~\emph{non-adaptively} chosen $\bfx$, and for lower-triangular $\bfC$'s when $\bfx$ is~\emph{adaptively} chosen (which is the typical situation for machine learning tasks).  We then demonstrate that our method yields both asymptotic improvements and experimental improvements in machine learning.

%
% \vspace{-0.25em}  % SPACE
\mypar{Conditional composition (\cref{sec:conditioningtrick}, \cref{thm:conditioningtrick})} This is our main technical tool that gracefully handles dependence between queries in the rows of $\bfC\bfx$; this arises due to multiple participations in rows of $\bfx$ for purposes of correlating noise. Specifically, we enable arbitrary queries $\bfC_{[i,:]}\cdot \bfx$ conditioned on $\bfC_{[1:i-1,:]}\cdot \bfx + \bfz_{1:i-1}$.
Standard composition theorems~\citep{dwork2014algorithmic} only handle this via a \textit{pessimistic worst-case privacy guarantee} that holds with certainty for each query.
%
\cref{thm:conditioningtrick} relaxes this to holding with high probability (over the randomness of the algorithm)
leading to significantly better guarantees. This generalizes an idea previously used in \citep{LocalToCentral,balle2020privacy} to analyze privacy amplification by shuffling. 
We believe this theorem will be useful
for analyzing correlated noise mechanisms beyond those studied herein. 

\mypar{Matrix mechanism privacy amplification via \MMCC{} (\cref{sec:cptb})} We prove amplified privacy guarantees for the matrix mechanism with uniform sampling, using \cref{thm:conditioningtrick}, that are nearly-tight in the low-epsilon regime as $\epsilon\to0$.
We improve over~\citet{BandedMF} because we enable ``more randomness'' in sampling---instead of participating w.p. $bp$ in $n/b$ rounds records can participate w.p. $p$ in all $n$ rounds.\anote{Check for consistency between samples and records.}

% \vspace{-0.25em}  % SPACE
Recall we need to analyze the privacy of outputting $\bfC \bfx + \bfz$, where rows of $\bfx$ are chosen via uniform sampling. We use \Cref{thm:conditioning} to reduce $\bfC \bfx + \bfz$ to a series of \textit{mixture of Gaussians (MoG) mechanisms} for which we can use privacy loss distribution (PLD) accounting. \MMCC{} is formally stated in \cref{fig:cptb}. 

\vspace{-0.25em}  % SPACE
\mypar{Binary tree analysis (\cref{sec:ftrl})} Letting $\sigma_{\epsilon, \delta}$ be the noise required for the Gaussian mechanism to achieve to satisfy $(\epsilon, \delta)$-DP, the binary tree mechanism requires noise $ \sigma_{\epsilon, \delta} \cdot \sqrt{\log n + 1}$. Owing to the versatility of conditional composition, we show that with shuffling, the (non-adaptive) binary tree mechanism only needs noise $\sigma_{\epsilon, \delta} \cdot O\left(\min\{\sqrt{\log n}, \sqrt{\log \log (1/\delta)}\}\right)$. This is optimal given current amplification by shuffling results, which require $n = \Omega(\log 1/\delta)$, We believe this requirement is necessary, but if one could show the current amplification by shuffling results hold for any $\delta$ then our upper bound would improve to $\sigma_{\epsilon, \delta} \cdot O(1)$. To the best of our knowledge, this is the first amplification guarantee (of any kind) for the binary tree mechanism. 

% \vspace{-0.25em}  % SPACE
\mypar{Empirical improvements (\cref{sec:empirical})} For our empirical studies, we write a library implementing \MMCC{}, which we are currently working on open-sourcing. The analysis of MoG mechanisms included in this library has other uses, such as tighter privacy guarantees for DP-SGD with group-level DP or for linear losses, see \cref{sec:implementation} for more discussion.

Using this library, first we show that $\epsilon$ computed via \MMCC{} for the binary tree mechanism matches the theoretical predictions of $\Omega(\sqrt{\log n})$ from \cref{sec:ftrl}. Then we apply our work to machine learning and show we can improve the privacy-utility tradeoff for binary-tree-DP-FTRL \citep{kairouz2021practical} entirely post-hoc. Finally, we empirically show that for the problem of minimizing $\ell_2^2$-error of all prefix-sums, a matrix mechanism analyzed with \MMCC{} gets smaller error than independent noise mechanisms for much smaller $\epsilon$ than past work. 

\subsection{Problem Definition}
\label{sec:probdef}

\mypar{Matrix mechanism \mm} Consider a workload matrix $\bfA\in\calR^{n\times n}$, and consider a data set $D=\{d_1,\ldots,d_m\}\in\calD^m$. Let $\bfx=\left[\bfx_1(D)|\cdots|\bfx_n(D)\right]^\top\in\calR^{n\times d}$ be a matrix s.t. each row $\bfx_i:\calD^*\to\calR^{d}$ is a~\emph{randomized} function that first selects a subset of the data set $D$ and then maps it to a real valued vector.
Further, each of the $\bfx_i$ has the following two properties. a)~\emph{Decomposability}: For the subset of the data set $D$ that $\bfx_i$ chooses (call it $S_i$), we have $\bfx_i(D)=\sum_{d\in S_i}\bfg_i(d)$ with $\bfg_i:\calD\to\calR^{d}$ is a vector valued function,
and b)~\emph{bounded sensitivity}: $\forall d\in\calD: \ltwo{\bfg_i(d)}\leq 1$. Observe that if this randomized function is also a) computing the flattened model gradient, b) clipping each per-example gradient, and c) averages the result, then this retrieves DP machine learning.

The class of (DP) \mm are those that approximate $\bfA\bfx$ with low-error (by minimizing some function of $\bfB\bfz$). Typically, one designs a pair of matrices $\bfB$ the decoder and $\bfC$ the encoder such that $\bfA=\bfB\bfC$ and $\bfC\bfx+\bfz$ satisfies DP\footnote{We use the zero-out adjacency \citep{ponomareva2023dpfy} to define DP in this paper.}, with $\bfz$ isotropic Gaussian noise. We assume $\bfC$ is non-negative for simplicity.

\mypar{Privacy amplification for the \mm} In this work we study the problem of amplifying the DP guarantee of the \mm if we incorporate the randomness in how the records of each $\bfx_i$ are selected (from $D$), e.g., how the minibatch is sampled.
We consider two selection strategies: 1)~\emph{uniform sampling}: each $\bfx_i$ selects each entry of $D$ independently w.p. $p$, and 2)~\emph{shuffling}: First the records of $D$ are randomly permuted, and then each $\bfx_i$ picks a fixed disjoint subset (of equal size) from $D$.

\mypar{Adaptivity} In our work we allow the choice of $\bfx_i$'s to be adaptive, i.e., $\bfx_i$ can be chosen based on the first $i-1$ outputs of \texttt{MM}. Under adaptivity, we will only consider encoder ($\bfB$) and decoder matrices ($\bfC$) that are lower triangular. However, for non-adaptive choices of the $\bfx_i$'s we  allow arbitrary choice of the matrices $\bfB$ and $\bfC$.~\emph{Unless mentioned specifically, all our results will be for the adaptive setting} as this pertains to ML.
